% ============================================================
%  Administrative Philosophy - Ronald J. Tackett, Ph.D.
%  Engine: XeLaTeX
%
%  CHANGELOG
%  v1.1  2026-09-04  SoFS start corrected from June to July 2025.
%  v1.0  2026-09-03  Rebuilt in LaTeX from the March 2026 PDF.
%                    Content unchanged; footer date moved to
%                    September 3, 2026 and the closing block from
%                    Spring 2026 to Fall 2026.
% ============================================================
\documentclass[10pt]{article}
\usepackage{tackettdoc}

\setdoclabel{Administrative Philosophy}
\setupdated{September 3, 2026}

\begin{document}

\statementheader
  {Founding Head, School of Foundational Studies}
  {Associate Professor of Physics \textperiodcentered\ Kettering University}
  {Administrative Philosophy}

\lede{Leadership is not a destination. It is a practice. And like any practice worth taking
seriously, it demands reflection, humility, and the willingness to keep learning, especially
when you thought you already knew.}

\stsection{I. Leadership Is Learned in Motion}

I became Founding Head of the School of Foundational Studies in July 2025, which means I
entered a formal administrative role at a moment of institutional complexity, with a unit that
did not yet exist and a community still finding its footing. I did not arrive with a finished
playbook. I arrived with a set of values and the understanding that the work itself would teach
me things no amount of preparation could.

That has proven true. Leadership does not arrive fully formed when a title changes. It
develops through decisions made with imperfect information, through conversations that do not
go as planned, and through the slow accumulation of judgment that comes only from being in the
room when difficult things happen. The ideals I started with, fairness, transparency,
advocating for my people, grounding decisions in principle, still matter. What has changed is
my understanding of how genuinely hard it is to live those ideals out consistently in real
institutions with real people.

I am still working out what good leadership looks like in practice, and I expect that to remain
true for as long as I am doing this work. That is not a disclaimer. It is the most honest thing
I can say about where I am.

\stsection{II. Authority and Trust Are Not the Same Thing}

One of the first things leadership taught me is that authority comes with a role, but trust is
earned through behavior over time. These are not the same thing, and confusing them is a
reliable way to alienate the people you are trying to lead.

People watch not just what leaders decide, but how those decisions are made. Who is consulted.
Who is surprised. Who feels heard, even when their preferred outcome does not prevail. Silence
from a leader is rarely neutral; it tends to be filled in with assumptions, and those
assumptions are rarely charitable. I have tried to learn from that, to communicate more than
feels necessary, to name things that might otherwise be inferred, and to acknowledge
uncertainty rather than projecting a confidence I do not always have.

Credibility, I have found, is less about projecting authority and more about being consistent,
accessible, and sincere. Those qualities accumulate slowly. They also erode quickly when
leaders communicate selectively or defensively. I try to keep that asymmetry in mind.

\stsection{III. Trust as the Foundation of Everything Else}

I think about trust more than almost anything else in my administrative work, partly because I
believe it is the scaffolding that holds a community together, and partly because I inherited a
community in which trust had been strained. The school I lead was formed through institutional
realignment that, whatever its necessity, was experienced by many as sudden and top-down. That
kind of experience leaves marks.

Trust, when it has been broken, cannot be restored by declaring it restored. It cannot be
rushed, and it cannot be accomplished through efficiency alone. Attempts to move on quickly
often deepen the damage by signaling that the relational work is secondary to operational
goals. Rebuilding trust requires a different posture entirely: consistency over persuasion,
presence over performance, and a genuine willingness to hear the same concerns more than once
without treating that repetition as a problem.

There is something important to understand about trust that I have come to believe deeply.
Trust exists only if all parties agree that it exists. If one person says it is not there, and
we take that seriously, then it simply is not there. It cannot be argued into existence,
declared by authority, or inferred from good intentions. When trust has been broken, the first
act of repair is acknowledging that it has been broken and taking that claim seriously, not
explaining it away.

\stsection{IV. Presence over Efficiency}

Modern institutional life is organized around efficiency. Meetings are shorter, messages are
faster, and decisions are expected to move quickly from discussion to execution. Many of the
tools that enable this pace are genuinely useful. But tools are not neutral. They shape what we
come to value. When efficiency becomes the dominant value, relationships tend to be the first
casualty.

Community does not form through optimization. It forms through repeated acts of showing up,
often in ways that feel inefficient. Trust is built in the margins: in conversations that
wander, in disagreements that get worked through rather than bypassed, in shared experiences
that accumulate over time. These processes resist streamlining. They require patience,
attentiveness, and presence.

I try to be present in ways that matter, not just visible in ways that are convenient. That
means staying in conversations longer than the agenda requires, seeking out colleagues I have
not heard from recently, and creating space for the kinds of informal interaction that do not
appear on any calendar but do much of the real work of culture-building. Choosing presence over
convenience will always feel slightly inefficient. That is precisely why it matters.

\stsection{V. Creating Conditions, Not Fixing Problems}

I came to administration from a technical and academic background, which means my default
instinct was problem-solving. Identify the issue, analyze the system, propose a solution,
implement, move on. That approach works reasonably well for equations and experiments. It works
far less well for organizations composed of individuals with complex histories, diverse
incentives, and strong attachments to how things have always been done.

What appears to be inefficiency from one vantage point may be adaptation from another. What
appears to be resistance may actually be fatigue, or grief, or a reasonable concern about being
left behind. Leadership forced me to slow down, especially when I felt pressure to move
quickly. Momentum without alignment is fragile. Progress that does not account for human
context rarely lasts.

I have had to learn that my role is often not to fix the problem directly, but to create
conditions in which others can engage it productively. That requires letting go of the
satisfaction of direct solutions and developing a different kind of patience: patience for
process, for distributed ownership, and for the slower pace at which people become genuinely
willing to change.

\stsection{VI. Honest Decisions and Accepted Discomfort}

Early in my leadership, I wanted consensus. I wanted meetings to end with agreement and a sense
of shared direction. Over time, I have come to understand that leadership sometimes requires
making principled decisions that not everyone will welcome, and that the discomfort this
produces is not automatically a sign of failure. Sometimes it is a sign that real tradeoffs are
being confronted honestly.

The challenge is finding the balance between decisiveness and openness. Lean too far into
decisiveness and you risk rigidity. Lean too far into accommodation and you risk paralysis.
Leadership lives in the tension between those two tendencies. There is no formula for resolving
it, only judgment, reflection, and a willingness to revisit decisions when new information
warrants it.

I also try to be honest with myself about my blind spots. Leadership amplifies both strengths
and weaknesses. For me, this has meant confronting tendencies to over-explain, to assume shared
context that is not actually shared, and to underestimate how emotionally charged certain
decisions feel to others, even when they appear procedural to me. You do not get to opt out of
your blind spots because your intentions are good.

\stsection{VII. Stewardship and What I Am Still Learning}

The framing I find most useful for the work I do is stewardship. I am temporarily entrusted
with people, structures, and cultures that often predate me and will certainly outlast me. My
task is not to remake them in my image but to leave them more coherent, more humane, and more
capable than I found them.

The most satisfying moments in this work are not when I personally solve a problem but when I
watch others step forward with confidence and ownership. That requires letting go of credit,
resisting the urge to intervene prematurely, and trusting people to grow into responsibility.
The goal, over time, is to make myself less central, not more.

I still misjudge situations. I still replay conversations and think about what I could have said
differently. I still get things wrong. What has changed is how I interpret those moments. I no
longer read them as evidence that I do not belong in this role. I read them as evidence that I
am paying attention. That shift, modest as it sounds, has made it possible to keep going.

\signatureblock{Fall 2026}

\end{document}
